% ---------------------------------------------------------------------------
% Summary of the full-year re-evaluation of sec:ch4-conditional-adaptation.
% Numbers generated from tables/ch4_*.csv (scripts/analyze_ch4_tier1.py).
% ---------------------------------------------------------------------------

\paragraph{Re-evaluation protocol.} The policy and all comparators were re-run over the full
campaign of each site on a common population: every sample whose clear-sky irradiance is at
least \SI{1}{\watt\per\square\meter}, the same daylight filter used in
\cref{ch:collaborative-learning}, and the same source checkpoint as that chapter. This gives
362~days and \num{245538} samples at Folsom, 298~days and \num{179858} at Karpicko, and 213~days
and \num{123523} at Poznań. The horizon is $H=\SI{15}{\minute}$ with
$T_1=\SI{10}{\watt\per\square\meter}$ and $T_2=\SI{100}{\watt\per\square\meter}$, and the error
$L$ of \cref{eq:ch4-abs-error} is computed on the reported, non-negative forecast. Each buffered
sample is released for an update only once \SI{15}{\minute} have elapsed since its prediction, so
every update uses a label that had already been measured; the gradient is evaluated with the
parameters in force at release time. The policy and the deterministic comparators are
bit-identical on repetition, so a single run is reported for them; few-shot adaptation is averaged
over three seeds, whose standard deviation in MAE is at most \SI{0.75}{\watt\per\square\meter}. Equivalence
with sliding-window adaptation is tested on relative MAE with a \SI{90}{\percent} moving-block
bootstrap interval (\num{7}-day blocks, \num{2000} resamples) against a $\pm\SI{5}{\percent}$
margin; pairwise differences in squared loss use Diebold--Mariano tests with a Newey--West
variance (lag~6), Holm-adjusted over all 171 contrasts.

\paragraph{The full-model action had not been exercised before.} In the implementation behind
the earlier results, the high-error branch of \cref{eq:ch4-update-trigger} unfroze the last two
children of the encoder, which for the backbone used here are a pooling layer and an identity map
without parameters. The high-error action therefore trained the head only, and the evaluated policy
was in effect head-only. The re-evaluation executes the two actions as specified, with
$\mathcal{S}_t=\mathbb{R}^{P}$ for $L\ge T_2$, and additionally reports the head-only policy
obtained by setting $T_2=\infty$.

\begin{table}[!t]
\caption[Full-year comparison of adaptation strategies on the daylight population.]{Full-year
comparison of adaptation strategies on the daylight population, $H=\SI{15}{\minute}$. MAE and RMSE
in \si{\watt\per\square\meter}, forecast skill (FS) in \si{\percent} against Smart Persistence
(\cref{eq:ch2-forecast-skill}). Bold marks the best value per site and column. Method key
as in \cref{tab:ch4-solar-benchmark}; the three conditional rows differ only in what the
high-error action ($L\ge T_2$) trains.}
\label{tab:ch4-full-year-benchmark}
\centering
\small
\TableDense
\begin{tabular}{@{}l l S[table-format=3.2] S[table-format=3.2] S[table-format=-2.2]@{}}
\toprule
\thead{Site} & \thead{Method} & \multicolumn{1}{c}{\thead{MAE $\downarrow$}} & \multicolumn{1}{c}{\thead{RMSE $\downarrow$}} & \multicolumn{1}{c}{\thead{FS $\uparrow$}}\\
\midrule
\multirow{11}{*}{Folsom} & Smart Persistence & 35.95 & 86.10 & 0.00\\
 & Static model & 31.51 & 72.60 & 15.68\\
 & MEMO & 33.26 & 72.74 & 15.51\\
 & NOTE & 104.79 & 138.62 & -61.01\\
 & RLS & 31.31 & 71.94 & 16.44\\
 & Window & 31.50 & \textbf{69.48} & \textbf{19.30}\\
 & Few-shot 1d & 31.78 & 73.15 & 15.04\\
 & Few-shot 7d & 37.14 & 74.75 & 13.18\\
 & Conditional, full model & 41.68 & 77.21 & 10.32\\
 & Conditional, head only & \textbf{31.03} & 70.16 & 18.51\\
 & Conditional, layer3+4, enc.\ LR $10^{-6}$ & 31.16 & 69.54 & 19.23\\
\midrule
\multirow{11}{*}{Karpicko} & Smart Persistence & 67.19 & 142.46 & 0.00\\
 & Static model & 80.40 & 140.14 & 1.63\\
 & MEMO & 79.58 & 140.90 & 1.09\\
 & NOTE & 134.07 & 169.51 & -18.99\\
 & RLS & 75.84 & 134.54 & 5.56\\
 & Window & 66.76 & \textbf{124.77} & \textbf{12.42}\\
 & Few-shot 1d & 68.65 & 136.60 & 4.11\\
 & Few-shot 7d & 68.53 & 138.38 & 2.87\\
 & Conditional, full model & 68.54 & 128.06 & 10.11\\
 & Conditional, head only & 66.58 & 127.99 & 10.16\\
 & Conditional, layer3+4, enc.\ LR $10^{-6}$ & \textbf{65.90} & 124.92 & 12.32\\
\midrule
\multirow{11}{*}{Poznań} & Smart Persistence & 67.73 & 138.99 & 0.00\\
 & Static model & 69.56 & 124.08 & 10.73\\
 & MEMO & 68.00 & 122.48 & 11.88\\
 & NOTE & 96.66 & 135.68 & 2.38\\
 & RLS & 66.56 & 120.84 & 13.06\\
 & Window & 58.92 & \textbf{111.79} & \textbf{19.57}\\
 & Few-shot 1d & 69.20 & 124.45 & 10.46\\
 & Few-shot 7d & 63.49 & 121.50 & 12.58\\
 & Conditional, full model & 62.68 & 115.37 & 16.99\\
 & Conditional, head only & 59.65 & 114.84 & 17.37\\
 & Conditional, layer3+4, enc.\ LR $10^{-6}$ & \textbf{58.58} & 112.70 & 18.91\\
\bottomrule
\end{tabular}
\end{table}

\paragraph{The full-model action degrades the model.} Executed as specified, the policy fails the
equivalence criterion at Folsom (\SI{+32.3}{\percent} relative MAE against sliding-window
adaptation, interval $[+26.4, +40.2]$) and at Poznań ($+6.4$, $[+4.8, +8.4]$). At Folsom, the
source domain, it ends \SI{32}{\percent} worse than the frozen model it started from
(\num{41.68} against \SI{31.51}{\watt\per\square\meter}). The degradation is progressive rather
than immediate: its monthly error relative to the frozen model rises from 1.06 in January to 4.38
in September. This is the residual risk retained in \cref{sec:ch4-adaptation-protocol}, a
trajectory bounded at every step that nonetheless drifts over a long horizon. With AdamW each
update moves every trainable weight by roughly the learning rate, whatever the size of its
gradient, and the full-model action was taken \num{24654} times at Folsom without any constraint
tying the encoder to its source values. Freezing the encoder removes this failure: at Folsom the full-model policy has
\SI{34}{\percent} higher MAE than the head-only one (Holm-adjusted $p<10^{-12}$),
while the two do not differ significantly in squared loss at Karpicko and Poznań.

\paragraph{Head-only adaptation matches sliding-window adaptation on typical errors.} The
head-only policy is equivalent to sliding-window adaptation at all three sites, with relative MAE
of $-1.5$ ($[-2.0, -1.0]$), $-0.3$ ($[-0.9, +0.4]$) and $+1.2$ ($[+0.7, +1.8]$) per cent, and has
the lowest MAE of any method at Folsom. Sliding-window adaptation keeps a lead on large errors,
with forecast skill higher by 0.8, 2.3 and 2.2 points and significantly lower squared loss at every
site ($p=0.042$, $0.0003$, $<0.0001$).

\begin{table}[!t]
\caption[Variants of the high-error action.]{Variants of the high-error action ($L\ge T_2$). The
trigger, thresholds, seed and checkpoint are identical; only the parameters trained by that action
and the stabiliser applied to it differ. Params: trainable parameters of the action in millions.
MAE in \si{\watt\per\square\meter}, FS in \si{\percent}. Eq.: relative MAE against sliding-window
adaptation within $\pm\SI{5}{\percent}$ at \SI{90}{\percent} confidence.}
\label{tab:ch4-high-error-variants}
\centering
\small
\TableDense
\begin{adjustbox}{max width=\textwidth}
\begin{tabular}{@{}l S[table-format=2.1] S[table-format=3.2] S[table-format=2.2] c S[table-format=3.2] S[table-format=2.2] c S[table-format=3.2] S[table-format=2.2] c@{}}
\toprule
 & & \multicolumn{3}{c}{\thead{Folsom}} & \multicolumn{3}{c}{\thead{Karpicko}} & \multicolumn{3}{c}{\thead{Poznań}}\\
\cmidrule(lr){3-5}\cmidrule(lr){6-8}\cmidrule(lr){9-11}
\thead{High-error action} & {\thead{Params}} & {\thead{MAE}} & {\thead{FS}} & \thead{Eq.} & {\thead{MAE}} & {\thead{FS}} & \thead{Eq.} & {\thead{MAE}} & {\thead{FS}} & \thead{Eq.}\\
\midrule
Full model & 24.0 & 41.68 & 10.32 & $\times$ & 68.54 & 10.11 & \checkmark & 62.68 & 16.99 & $\times$\\
Layer3+4 & 22.6 & 37.43 & 15.42 & $\times$ & 66.00 & 13.46 & \checkmark & 59.96 & 19.89 & \checkmark\\
Layer3+4 + source replay & 22.6 & 37.05 & 13.71 & $\times$ & 66.32 & 14.05 & \checkmark & 60.92 & 19.37 & $\times$\\
Layer3+4 + decoupled L2-SP & 22.6 & 31.44 & 17.82 & \checkmark & 67.76 & 8.24 & \checkmark & 59.45 & 17.80 & \checkmark\\
Layer3+4, enc.\ LR $10^{-6}$ & 22.6 & 31.16 & 19.23 & \checkmark & 65.90 & 12.32 & \checkmark & 58.58 & 18.91 & \checkmark\\
Head only & 0.5 & 31.03 & 18.51 & \checkmark & 66.58 & 10.16 & \checkmark & 59.65 & 17.37 & \checkmark\\
\midrule
Sliding window (reference) & {--} & 31.50 & 19.30 & {--} & 66.76 & 12.42 & {--} & 58.92 & 19.57 & {--}\\
\bottomrule
\end{tabular}
\end{adjustbox}
\end{table}

\paragraph{Partial encoder updates and stabilisers.} Restricting the high-error action to the
last two residual stages (layer3 and layer4, \SI{94}{\percent} of the encoder) slows the drift
without stopping it: at Folsom the policy remains \SI{19}{\percent} worse than the frozen model.
At the shifted sites, however, it raises forecast skill above sliding-window adaptation (13.46 and
19.89 against 12.42 and 19.57 at Karpicko and Poznań), so updating the deep layers does reduce
large errors where there is shift to absorb.
Three stabilisers of that action were compared (\cref{tab:ch4-high-error-variants}): appending
seven source-domain training samples to each encoder step, a decoupled pull of each encoder weight
towards its source value after every step (\SI{0.5}{\percent} of the displacement), and a
tenfold smaller encoder learning rate. None both protects the source domain and preserves the
large-error gain at the shifted sites. Replay leaves the Folsom result essentially unchanged (\num{37.05} against
\SI{37.43}{\watt\per\square\meter}), consistent with it altering the direction of the encoder
steps but not their size. The decoupled pull protects Folsom
but prevents useful adaptation at Karpicko. The reduced encoder learning rate is the best
compromise. It has the lowest MAE of any method at Karpicko and Poznań and is within
\SI{0.13}{\watt\per\square\meter} of the head-only policy at Folsom, is significantly better than
head-only adaptation in squared loss at all three sites ($p=0.001$, $0.0002$, $<0.0001$), matches
sliding-window adaptation in squared loss at Folsom and Karpicko ($p=1.0$), and trails it at
Poznań ($p=0.030$).

\paragraph{Update rate.} Error-triggered updates are frequent. The head-only policy updates on
\SI{42.7}{\percent}, \SI{69.0}{\percent} and \SI{65.2}{\percent} of released samples at Folsom,
Karpicko and Poznań; in the variants that reach the encoder, the high-error action accounts for
\SIrange{8.6}{10.3}{\percent} of decisions at Folsom and \SIrange{17.0}{19.8}{\percent} at the
Polish sites.

\paragraph{Limitations.} The comparison rests on a single source checkpoint; sensitivity to it
is not characterised. Whether the error threshold itself contributes, as opposed to the update rate
it produces, is not tested, since no rate-matched random trigger was run. The execution benchmark
of \cref{tab:ch4-solar-edge-benchmark} was measured with the earlier implementation, so it
characterises head-only behaviour and does not reflect the cost of encoder updates.
